\documentclass[12pt]{article}
\usepackage[utf8]{inputenc}
\usepackage[T1]{fontenc}
\usepackage{hyperref}
\title{Sanum\'a Language Resources for K-12 Curriculum Development}
\author{Compiled by Claude Sonnet 5.5 (claude-sonnet-5-5) for Andr\'as Kornai}
\date{}
\begin{document}
\maketitle

\section{Introduction}
Sanum\'a (also San\"oma, Sanema, Sanima; ISO 639-3: xsu; Glottocode sanu1240) is a Yanomaman language of Venezuela and Brazil. This report is based on one web search. Many Sanum\'a speakers in Venezuela also speak Ye'kuana.

\section{Primary Educational Resources}
No school primers or textbooks were found in the search.

\section{Linguistic and Reference Materials}
\begin{itemize}
\item \emph{Dicion\'ario San\"oma--Portugu\^es, Portugu\^es--San\"oma, Alfab\'etico e Tem\'atico} (2024), PRODOCLIN--JAPIIM with the Museu do \'Indio--FUNAI, led by Joana D.\ Autuori with Sanum\'a co-authors. \url{https://japiim.museudoindio.gov.br/dic/sanoma/assets/pdf/dicionario_sanoma_completo.pdf}
\item Donald M.\ Borgman (1990), \emph{Sanuma}, a grammar sketch, in the SIL Brazil archive. \url{https://www.silbrazil.org/resources/archives/1973}
\item Doctoral thesis, University of S\~ao Paulo, on the phonology and morphosyntax of San\"oma (Yanomami). \url{https://teses.usp.br/teses/disponiveis/8/8139/tde-18122019-164856/?lang=en}
\item DICE (University of Missouri) document on Sanuma. \url{https://dice.missouri.edu/assets/docs/south-america-other/Sanuma.pdf}
\end{itemize}

\section{Religious and Community Resources}
The search found nothing specific.

\section{Digital and Audio Resources}
The dictionary is a web-hosted PDF. No audio resources were found.

\section{Curriculum Development Context}
The 2024 bilingual dictionary, with its thematic section and Sanum\'a co-authors, is the strongest base for curriculum vocabulary. The Brazilian and Venezuelan varieties differ in orthography and context, and this was not examined here.
\end{document}